\section{Задача}
\begin{definition}[Вариант 34: сортировочный центр посылок]
Посылка с уникальным номером проходит входной накопитель, сканер, сортировщик, конвейер и выход своего направления. Вместимости очередей и производительность устройств ограничены. При заполнении следующего накопителя устройство ждёт с посылкой. Ошибка чтения вызывает повтор; после исчерпания попыток посылка поступает в ручной накопитель. Его дальнейшая обработка не моделируется.
\end{definition}

\section{Потоки и синхронизация}
Реализация использует C11, libc, системные вызовы и pthread.
\begin{center}
\begin{tabular}{p{0.27\linewidth}p{0.09\linewidth}p{0.51\linewidth}}
\toprule
Участник & Число & Поведение \\
\midrule
Посылки & $N$ & Ожидание результата, повтор, переход \\
Входные линии & $L$ & Подача с переменными интервалами \\
Сканеры, сортировщики & $S+R$ & Чтение маркировки, сортировка \\
Конвейеры & $D$ & Перевозка за $\lceil l/v\rceil$ единиц времени \\
Разгрузчики & $D+1$ & Освобождение выходов и ручного накопителя \\
Служебные & $2$ & Главный поток и приём сигналов \\
\bottomrule
\end{tabular}
\end{center}
Всего $N+L+S+R+2D+3$ потока; по умолчанию 35 при 20 посылках.

\begin{remark}[Согласование работы]
Кольцевые очереди используют семафоры свободных мест и элементов. Общий мьютекс защищает очереди, состояния и журнал; ожидания и задержки выполняются без него. Устройство сообщает посылке результат через семафор. Посылка получает место в следующей очереди, переходит и подтверждает освобождение устройства. Повтор выполняется на том же сканере.
\end{remark}

\section{Корректность и завершение}
\texttt{-{}-check} проверяет единственное местонахождение посылки, уникальность номера, вместимости, занятость устройств и распознавание направления перед сортировкой. Итоговая проверка выполняется всегда. При прерывании сохраняется баланс доставленных, отказов, незавершённых и ещё не поступивших посылок.

После доставки или отказа всех посылок программа возвращает 0. Сигналы SIGINT и SIGTERM принимает поток с \texttt{sigwait}: пробуждает ожидающих, задержки проверяют остановку каждые 20 мс. После \texttt{pthread\_join} сохраняются итоги и освобождаются ресурсы. Коды: 130/143 при сигналах, 2 при неверных аргументах, 1 при ошибке выполнения.

\clearpage
\section{Интерфейс и запуск}
Аргументы задают размеры партии и очередей, число устройств, задержки и ошибки; справка - \texttt{-{}-help}. Допускаются 0--512 посылок и 1--64 устройства каждого вида.
\begin{lstlisting}[frame=none,backgroundcolor=\color{example_back},aboveskip=0pt,belowskip=0pt,xleftmargin=0pt,xrightmargin=0pt]
make
./build/sorting_center --parcels 17 --error-percent 100 \
    --retries 3 --check --log run.log
make test
\end{lstlisting}
\texttt{-{}-delay-ms}: 1 мс по умолчанию, 0 отключает задержки. Seed определяет результаты посылок, но не порядок событий. Консоль и \texttt{-{}-log} показывают время, номер посылки, переходы, ожидания, ошибки и итоги. Без аргументов используются стандартные параметры.

\section{Проверка и результаты}
На macOS с Apple Clang 17 пройдены 14 групп тестов: маршруты, FIFO, вместимости, параллельность и длительность обработки. Проверены 25 случайных конфигураций, ёмкости 1, 3 и 7 и десять повторных запусков.

\begin{center}
\begin{tabular}{p{0.51\linewidth}p{0.40\linewidth}}
\toprule
Сценарий & Результат \\
\midrule
40 посылок, ошибок нет & 40 доставок, 40 чтений \\
17 посылок, 100\% ошибок, 3 повтора & 17 отказов, 68 чтений \\
20 посылок, ёмкость 1 & 20 доставок; есть ожидание \\
Пустая партия & Код 0, нулевые счётчики \\
SIGINT/SIGTERM: задержка, полный выход & Коды 130/143, баланс сохранён \\
Отказ создания потоков & Код 1, потоки завершены \\
\bottomrule
\end{tabular}
\end{center}
Проверены неверные аргументы, ошибки вывода и шесть точек отказа создания потоков. Проверка под Linux и санитайзерами не завершена; запуск в Linux - \path{tests/linux.sh}.

\section{Ограничения и использование ИИ}
\begin{exercise}[Критическая оценка]
Очереди работают за $O(1)$, полная проверка события - за $O(N+K)$. Данные занимают $O(N+B+K)$ плюс стеки потоков; $K$ - число устройств и очередей, $B$ - их суммарная вместимость. Общий мьютекс и медленный вывод ограничивают скорость, отдельный поток каждой посылки - размер партии.

Исправлен разбор времени в тестовом журнале; сборки с санитайзерами разнесены по каталогам. При отказе создания потоков уже запущенные останавливаются и присоединяются.
\end{exercise}

ИИ применялся для модели, тестов и отчёта; примеры запросов - в \path{docs/PROMPTS.md}. Из первого ИДЗ сохранены разбор аргументов, журнал и оформление. Реализована одна версия с мьютексом и семафорами.
